\section{Finite obstruction bases}
\label{sec:finite-obstructions-universal}

The two-state case proof supplies three levels of obstruction uniformity:
\begin{equation}
\label{eq:three-obstruction-quantifiers}
\begin{aligned}
  &\text{controller-dependent trap complexity:}
  &&\forall A\;\exists I_A,\\
  &\text{finite complete test family:}
  &&\exists\mathcal U\text{ finite}\;\forall A\;\exists I\in\mathcal U,\\
  &\text{one universal host maze:}
  &&\exists U\;\forall A\;\exists(x_A,y_A)\in U^2.
\end{aligned}
\end{equation}
Here $A$ ranges over legal controllers with at most two states, and each test
or start pair is initially noncontacting. The successive lines fix the controller,
then a menu, then the host itself. Finite controller space already gives an abstract
finite menu from $F(2)<\infty$; the construction below instead extracts a small
explicit family from the proof without enumerating controllers.

\subsection{An explicit proof-derived obstruction basis}
\label{subsec:human-obstruction-basis}

A \emph{test instance} is a triple $I=(V,x_0,y_0)$ consisting of an absolutely
oriented finite connected induced grid maze and two distinct state-$0$ starts with
$\|x_0-y_0\|_1\ge2$.  We say that $I$ \emph{traps} $A$ when the two synchronous
copies of $A$ started at $x_0,y_0$ never reach Manhattan distance at most one.
Compass orientation is part of the instance.

\begin{theorem}[Explicit finite obstruction basis]
\label{thm:human-finite-basis}
There is an explicit family $\mathcal U_{\mathrm{hum}}$ of $144$ absolutely oriented,
start-labelled test instances, each on at most seven cells, such that every legal
deterministic controller with at most two states is trapped by at least one member of
$\mathcal U_{\mathrm{hum}}$.

The family uses $76$ fixed oriented maze geometries.  Under simultaneous $D_4$
transport and exchange of the two identical agents it has $22$ instance orbits, and
its maze geometries have $13$ $D_4$-types.
\end{theorem}

\begin{proof}
The upper proof has $30$ primitive terminal predicates.  Making explicit the
two equal-direction variants and the four reset-bit variants produces $34$ normalized
source records.  In every branch the proof applies one global element $g\in D_4$ to
the controller, masks, directions, maze and starts; it never rotates only the maze.
The selected normalized trap is therefore transported back by $g^{-1}$.  We take all
eight inverse transports of all $34$ source records, then quotient only by common
translation, exchange of the two identical state-$0$ agents, and literal identity.
Rotation and reflection are \emph{not} quotient operations on the concrete family,
because the compass is absolute.  This yields exactly the $144$ instances stated
above.

The completeness implication itself is the original case tree.  For an
arbitrary two-state controller that tree reaches one of its terminal source records;
the corresponding normalized execution is trapped by the terminal lemma.  Simultaneous
$D_4$ conjugacy preserves tagged executions and Manhattan distance, so the inverse
transport traps the original controller.  That transported instance belongs to
$\mathcal U_{\mathrm{hum}}$.  A one-state controller is included by viewing it as a
two-state table whose state-$0$ rows always return state $0$.
\end{proof}

The counting in Theorem~\ref{thm:human-finite-basis} is finite bookkeeping rather
than a premise of completeness.  The $34$ normalized source records and all $144$
transported instances are supplied in the supplementary data.  Their terminal
conditions organize into the eight proof modules summarized in
Table~\ref{tab:obstruction-modules}.

% Exact selector summary. Full terminal proofs stay in section 5 and appendix A.
\begin{table}[!htbp]
\centering\small
\setlength{\tabcolsep}{4pt}
\renewcommand{\arraystretch}{1.08}
\begin{tabularx}{\linewidth}{@{}lP{.44\linewidth}Y@{}}
\toprule
Module & Regime / first failed condition & Geometry and noncontact mechanism \\
\midrule
$\mathrm{AX}$ & Straight directions coincide, or a normalized opposed-direction toggle/reset/endpoint test fails.
& Line or U-path, $\le7$ cells. Terminal shuttles or disjoint line supports; parity covers adjacent supports. \\
\addlinespace[3pt]
$\mathrm{RP}$ & After the front tests, both axes reset: $C_i,C_j\in\{C_0,C_1\}$.
& Two three-edge arms, $7$ cells. Closed pockets leave the shared corner unvisited. \\
\addlinespace[3pt]
$\mathrm Z$ & One axis is $C_0$, the other $I$; exchange axes if needed. Six successive corner alternatives.
& L/U/offset path, $\le7$ cells. Closed arm regions or forced prefixes; separation or parity. \\
\addlinespace[3pt]
$\mathrm O$ & One axis is $C_1$, the other $I$. Seven successive corner alternatives.
& L/U/offset path, $\le7$ cells. Closed pockets and forced corner prefixes; separation or parity. \\
\addlinespace[3pt]
$\mathrm{LOCK}$ & Both axes are $I$. After drift/diagonal normalization, the first of three corner-direction tests fails.
& Paired terminal-arm path, $\le7$ cells. Locked disjoint translated supports. \\
\addlinespace[3pt]
$\mathrm{BIT}$ & All three direction tests pass; the first of four required corner output bits fails.
& Six-cell path. Closed tagged regions or a prefix into disjoint cycles; gap or parity. \\
\addlinespace[3pt]
$\mathrm{TAIL}$ & Lock and bit tests pass, but the final two W9 entries do not both hold.
& Seven-cell path. A forced shuttle versus a separated left region, including the transient. \\
\addlinespace[3pt]
$\mathrm J$ & W9 holds. For $u=\delta(0,NSW)$, $v=\delta(1,NSW)$, choose G2 if $v=(S,0)$ or $[u=(S,0),\ v\in\{(N,0),(W,1)\}]$; otherwise G1.
& Seven-cell one-junction tree. Exact tagged words separate the executions through the transient. \\
\bottomrule
\end{tabularx}
\caption{The eight terminal modules: $C_0,C_1$ are reset next-state maps and $I$ is identity. These selectors guide the complete terminal proofs; only $\mathrm J$ needs a junction rather than a path.}
\label{tab:obstruction-modules}
\end{table}


\subsection{A computer-minimized subfamily}
\label{subsec:44-subfamily}

The proof-derived family is intentionally faithful to the proof rather than optimized for
cardinality.  Exact computation gives a smaller complete menu, but its scope must be
stated carefully.

\begin{proposition}[Minimum within the proof-derived candidate family]
\label{prop:44-candidate-minimum}
There exists a complete test family of $44$ instances, all chosen from
$\mathcal U_{\mathrm{hum}}$.  Moreover, no complete subfamily of
$\mathcal U_{\mathrm{hum}}$ has fewer than $44$ members.
\end{proposition}

The existence and candidate-restricted optimum in this proposition are
computer-assisted results, unlike the deductive completeness of the full family.
\paragraph{Certificate sketch.}
For the upper bound, the full physical transition function of each of the $144$
instances was reconstructed, and a SAT encoding of a legal controller escaping the
selected $44$ instances was proved unsatisfiable by an independently checked LRAT
certificate.  For the lower bound, there are $44$ explicit legal controllers
$P_1,\ldots,P_{44}$.  If
\[
  S_i=\{I\in\mathcal U_{\mathrm{hum}}: I\text{ traps }P_i\},
\]
then each $S_i$ is nonempty and the $44$ sets are pairwise disjoint.  Any complete
subfamily of the $144$ candidates must therefore choose a distinct member from each
$S_i$.  This packing argument is a direct exact finite check of the full physical
instances and does not use SAT.

The $44$ tests are complete, but optimality is \emph{restricted to the $144$
supplied candidates}. Whether an arbitrary complete family can have at most $43$
members remains open.
